\section{结论}

以“算法清单”为起点，模型、代码和写作链才是本轮可核对的产物。三名成员分别
承担机理连续模型、综合优化和数据决策的首责，编程、误差分析和选题判断则共同承担；
18 道真题分别提供练习入口。论文研究以 59 篇编号语料为限，59 张可追溯的正式证据卡覆盖
2892 页；赛题和专家评述不混入文风样本，A/B/C 三类分别完成 20/19/20 篇人工门禁。
模板与 Skill 可随本轮 59 篇证据增量修改；四周计划把尚无实证的能力落实为代码、图表、
论文和复现产物，并逐项核对。后续限时演练不增加算法名称，重点检验基线建立、约束书写、
关键结果复算和文字说明能否在时限内完成。

\begin{thebibliography}{99}
\bibitem{papers}
2020--2025 年全国大学生数学建模竞赛 A/B/C 题编号论文语料（本地归档），
59 篇，2892 页。

\bibitem{problems}
全国大学生数学建模竞赛赛题原文：2020--2025 年 A/B/C 题（本地归档）。

\bibitem{commentaries}
全国大学生数学建模竞赛专家评述：2020--2025 年相关题目（本地归档）。

\bibitem{mindmap}
数学建模学习思维导图 \texttt{1.png}（用户提供，本地文件，核验日期：
2026-07-28）。
\end{thebibliography}

\appendix
